% ============ 第七章 §4 分歧指数和剩余类域次数 (原书 pp.261-266) ============
\section{分歧指数和剩余类域次数}\label{sec:ch7-s4}

以上研究了同一个域中不同赋值的分类（素除子）. 本节研究域 $E$ 和它的子域 $F$ 的赋值和素除子之间的联系. 当 $E/F$ 是有限（代数）扩张, 并且均是非阿赋值域时, 我们要研究 $E$ 和 $F$ 在赋值结构上的联系.

设 $F$ 是域 $E$ 的子域, $\varphi$ 和 $\psi$ 分别是域 $E$ 和 $F$ 的赋值, 如果 $\psi=\varphi|_F$, 我们称 $\psi$ 是 $\varphi$ 在 $F$ 的\emph{限制}, 而 $\varphi$ 叫作 $\psi$ 到 $E$ 的\emph{扩充}, 并表示成 $(F,\psi)\subseteq(E,\varphi)$, 叫作赋值域的扩张. 由第 $7.2$ 节可知, $\varphi$ 和 $\psi$ 或者同时为非阿的, 或者同时是阿基米德的. 我们还可证明:

\begin{lemma}\label{lem:ch7-8}
设 $(F,\psi)\subseteq(E,\varphi)$ 是赋值域的代数扩张, 则: $\psi$ 为平凡赋值当且仅当 $\varphi$ 为平凡赋值.
\end{lemma}

\begin{proof}
若 $\varphi$ 平凡, 则 $\psi$ 显然平凡. 现在设 $\psi$ 是平凡赋值, 即对每个 $a\in F^{\times}$, $\psi(a)=1$, 则 $\psi$ 和 $\varphi$ 均是非阿的. 我们用反证法证明 $\varphi$ 为 $E$ 的平凡赋值. 若不然, 则存在 $\alpha\in E^{\times}$, 使得 $\varphi(\alpha)>1$. 由于 $E/F$ 是代数扩张, 从而有
\[
  \alpha^n+a_1\alpha^{n-1}+\cdots+a_n=0,\qquad a_i\in F,\ 1\leqslant i\leqslant n,\ n\geqslant 1.
\]
但是 $\varphi(a_i)=\psi(a_i)=0$ 或 $1$, 可知 $\varphi(\alpha^n)$ 比上式左边其余项的赋值都大, 这就与引理 \ref{lem:ch7-5} 相矛盾.
\end{proof}

以下设 $E/F$ 是域的代数扩张. 设 $\varphi_1$ 和 $\varphi_2$ 为 $E$ 的非平凡赋值. 由上面引理知它们在 $F$ 的限制 $\psi_1$ 和 $\psi_2$ 也是非平凡的. 如果 $\varphi_1$ 和 $\varphi_2$ 等价（$\varphi_2=\varphi_1^{\alpha}$）, 则 $\psi_1$ 和 $\psi_2$ 也等价（$\psi_2=\psi_1^{\alpha}$）, 从而 $E$ 的一个素除子 $Q$ 中所有赋值在 $F$ 的限制属于 $F$ 的同一个素除子 $P$. 我们称 $P$ 为 $Q$ 的\emph{限制}, $Q$ 为 $P$ 的\emph{扩充}, 表示成 $P=Q|_F$, $(F,P)\subseteq(E,Q)$, 而且 $P$ 和 $Q$ 或者均为非阿素除子, 或者均是阿基米德素除子.

本节以下设 $(F,P)\subseteq(E,Q)$ 是非阿赋值域的有限扩张（从而 $E/F$ 是代数扩张）. 以 $O_E$, $\mathfrak{p}_E$ 表示 $(E,Q)$ 的赋值环和素理想, $O_F$ 和 $\mathfrak{p}_F$ 是 $(F,P)$ 的赋值环和素理想, 则 $\overline{E}=O_E/\mathfrak{p}_E$ 和 $\overline{F}=O_F/\mathfrak{p}_F$ 分别是 $(E,Q)$ 和 $(F,P)$ 的剩余类域. 由 $O_F\subseteq O_E$ 和 $\mathfrak{p}_F=\mathfrak{p}_E\cap O_F$, 可知 $\overline{F}$ 可看成 $\overline{E}$ 的子域. 下面引理 \ref{lem:ch7-9} 要证 $\overline{E}/\overline{F}$ 是有限扩张. 另一方面, 设 $v$ 是 $Q$ 中一个指数赋值, 则 $v(F^{\times})$ 为 $v(E^{\times})$ 的加法子群. 易知 $[v(E^{\times}):v(F^{\times})]$ 和 $v$ 在 $Q$ 中选取方式无关, 我们也要证明 $[v(E^{\times}):v(F^{\times})]$ 是有限的.

\begin{definition}\label{def:ch7-6}
设 $(F,P)\subseteq(E,Q)$ 是非阿赋值域的有限扩张, $v$ 为 $Q$ 中一个指数赋值. 称
\[
  f=f(Q/P)=[\overline{E}:\overline{F}]
\]
为 $Q$ 对 $P$ 的\emph{剩余类域次数}, 而
\[
  e=e(Q/P)=[v(E^{\times}):v(F^{\times})]
\]
为 $Q$ 对 $P$ 的\emph{分歧指数}.
\end{definition}

由定义易知, 若 $(F,P)\subseteq(E,Q)\subseteq(K,R)$ 是非阿赋值域的有限扩张, 则
\[
  e(R/P)=e(R/Q)e(Q/P),\qquad f(R/P)=f(R/Q)f(Q/P).
\]

\begin{lemma}\label{lem:ch7-9}
设 $(F,P)\subseteq(E,Q)$ 为非阿赋值域的有限扩张, $v$ 为 $Q$ 中一个指数赋值, 则

(1) 如果 $w_1,\cdots,w_r\in O_E$, 并且 $\overline{E}$ 中元素 $\bar{w}_1,\cdots,\bar{w}_r$ 是 $\overline{F}$ 线性无关的, 则对于 $a_1,\cdots,a_r\in F$, 必然有
\[
  v(a_1w_1+\cdots+a_rw_r)=\min_{1\leqslant i\leqslant r}\{v(a_i)\}.
\]
特别地, $w_1,\cdots,w_r$ 是 $F$-线性无关的. 于是 $f(Q/P)\leqslant[E:F]$.

(2) 如果 $\pi_0,\cdots,\pi_s\in E^{\times}$, 并且 $v(\pi_j)\ (0\leqslant j\leqslant s)$ 是 $v(E^{\times})$ 对于 $v(F^{\times})$ 的陪集代表元素, 则对 $b_0,\cdots,b_s\in F$,
\[
  v(b_0\pi_0+\cdots+b_s\pi_s)=\min_{0\leqslant i\leqslant s}\{v(b_i\pi_i)\}.
\]
特别地, $\pi_0,\cdots,\pi_s$ 是 $F$-线性无关的, 于是 $e(Q/P)\leqslant[E:F]$.
\end{lemma}

\begin{proof}
(1) 不妨设 $a_i\ (1\leqslant i\leqslant r)$ 不全为零, 并且不妨设 $v(a_1)=\min\limits_{1\leqslant i\leqslant r}\{v(a_i)\}$, 则 $v(a_1)<\infty$, 并且 $a_i/a_1\in O_F\ (1\leqslant i\leqslant r)$. 从而
\[
  \alpha=w_1+\frac{a_2}{a_1}w_2+\cdots+\frac{a_r}{a_1}w_r\in O_E.
\]
由于 $\bar{w}_1,\cdots,\bar{w}_r$ 是 $\overline{F}$-线性无关的, 可知
\[
  \bar{\alpha}=\bar{w}_1+\overline{\left(\frac{a_2}{a_1}\right)}\bar{w}_2+\cdots+\overline{\left(\frac{a_r}{a_1}\right)}\bar{w}_r\neq 0\in\overline{E}.
\]
于是 $\alpha\in O_E-\mathfrak{p}_E$, 即 $v(\alpha)=0$, 从而
\begin{align*}
  v(a_1w_1+\cdots+a_rw_r)&=v(a_1)+v(\alpha)=v(a_1)\\
  &=\min_{1\leqslant i\leqslant r}\{v(a_i)\}.
\end{align*}
由此即知 $f(Q/P)=[\overline{E}:\overline{F}]\leqslant[E:F]$.

(2) 令 $\beta=b_0\pi_0+\cdots+b_s\pi_s$, 先设 $b_j\ (0\leqslant j\leqslant s)$ 均不为 $0$, 则 $v(b_j\pi_j)<\infty\ (0\leqslant j\leqslant s)$, 我们证明 $v(b_j\pi_j)\ (0\leqslant j\leqslant s)$ 彼此不同. 因若 $i\neq j$ 而 $v(b_i\pi_i)=v(b_j\pi_j)$, 则
\[
  v(\pi_i)-v(\pi_j)=v(\pi_i/\pi_j)=v(b_j/b_i)\in v(F^{\times}),
\]
这与 $v(\pi_i)$ 和 $v(\pi_j)$ 属于 $v(E^{\times})$ 对 $v(F^{\times})$ 的不同陪集相矛盾. 于是 $v(\beta)=\min\limits_{0\leqslant j\leqslant s}\{v(b_j\pi_j)\}$. 若某些 $b_j$ 为 $0$, 则 $\beta=\sum_{b_j\neq 0}b_j\pi_j$. 而由上述证明知
\[
  v(\beta)=\min_{b_j\neq 0}\{v(b_j\pi_j)\}=\min_{0\leqslant j\leqslant s}\{v(b_j\pi_j)\}.
\]
\end{proof}

上面引理可以改进为

\begin{lemma}\label{lem:ch7-10}
设 $(E,Q)\supseteq(F,P)$ 是非阿赋值域的有限扩张, 则 $e(Q/P)f(Q/P)\leqslant[E:F]$.
\end{lemma}

\begin{proof}
取 $v$ 为 $Q$ 中一个指数赋值, 记
\begin{align*}
  e&=e(Q/P)=[v(E^{\times}):v(F^{\times})],\\
  f&=f(Q/P)=[\overline{E}:\overline{F}],
\end{align*}
则存在 $w_1,\cdots,w_f\in O_E$, 使得 $\bar{w}_1,\cdots,\bar{w}_f$ 是 $\overline{F}$-线性无关的, 又存在 $\pi_0,\cdots,\pi_{e-1}\in E^{\times}$, 使得 $v(\pi_j)\ (0\leqslant j\leqslant e-1)$ 为 $v(E^{\times})$ 对 $v(F^{\times})$ 的陪集完全代表系. 我们要证 $w_i\pi_j\ (1\leqslant i\leqslant f,\ 0\leqslant j\leqslant e-1)$ 是 $F$-线性无关的, 从而 $ef\leqslant[E:F]$. 为此, 我们只需证明: 对任意 $a_{ij}\in F\ (1\leqslant i\leqslant f,\ 0\leqslant j\leqslant e-1)$, 有
\[
  v\left(\sum_{i,j}a_{ij}w_i\pi_j\right)=\min_{i,j}\{v(a_{ij}\pi_j)\}.
\]
不妨设 $a_{ij}$ 不全为 $0$, 令
\[
  J=\{j\mid 0\leqslant j\leqslant e-1,\ \text{存在 } i,\ \text{使得 } a_{ij}\neq 0\},
\]
则 $J$ 不是空集, 并且
\[
  \sum_{i,j}a_{ij}w_i\pi_j=\sum_{j\in J}\pi_j\sum_ia_{ij}w_i.
\]
由引理 \ref{lem:ch7-9} 知当 $j\in J$ 时,
\[
  v\left(\sum_ia_{ij}w_i\pi_j\right)=v(\pi_j)+v\left(\sum_ia_{ij}w_i\right)
  =v(\pi_j)+\min_i\{v(a_{ij})\},
\]
于是对 $J$ 中不同的 $j$, 上式右边的值属于 $v(E^{\times})$ 对 $v(F^{\times})$ 的不同陪集, 从而彼此不同. 所以
\begin{align*}
  v\left(\sum_{ij}a_{ij}w_i\pi_j\right)&=v\left(\sum_{j\in J}\pi_j\sum_ia_{ij}w_i\right)\\
  &=\min_{j\in J}\left\{v(\pi_j)+\min_i\{v(a_{ij})\}\right\}\\
  &=\min_{i,j}\{v(a_{ij}\pi_j)\}.
\end{align*}
\end{proof}

再进一步改进则有

\begin{theorem}\label{thm:7.5}
设 $(E,Q_i)\supseteq(F,P)\ (1\leqslant i\leqslant g)$ 是非阿赋值域的有限扩张, 其中 $Q_1,\cdots,Q_g$ 为 $P$ 到 $E$ 的不同扩充素除子, 则
\[
  \sum_{i=1}^{g}e(Q_i/P)f(Q_i/P)\leqslant[E:F].
\]
\end{theorem}

\begin{proof}
不妨设 $P$ 非平凡, 从而 $Q_i$ 也非平凡. 令 $O$ 和 $\mathfrak{p}$ 是 $F$ 对 $P$ 的赋值环和素理想, $O_i$ 和 $\mathfrak{p}_i$ 为 $E$ 对 $Q_i$ 的赋值环和素理想（$1\leqslant i\leqslant g$）. 取 $\varphi\in P$, $\varphi^{(i)}\in Q_i$, 使得 $\varphi^{(i)}|_F=\varphi\ (1\leqslant i\leqslant g)$. 令
\[
  \psi\colon O\to\overline{F}=O/\mathfrak{p},\qquad
  \psi^{(i)}\colon O_i\to\overline{E}^{(i)}=O_i/\mathfrak{p}_i\ (1\leqslant i\leqslant g)
\]
均是环的自然满同态, 则 $\psi^{(i)}|_O=\psi$. 记
\[
  e_i=e(Q_i/P),\qquad f_i=f(Q_i/P)\qquad(1\leqslant i\leqslant g).
\]
我们像引理 \ref{lem:ch7-10} 中那样取
\[
  \{w_{\nu}^{(i)}\in O_i\mid 1\leqslant\nu\leqslant f_i\},\qquad
  \{\pi_{\mu}^{(i)}\in E^{\times}\mid 0\leqslant\mu\leqslant e_i-1\}
\]
使得 $\psi^{(i)}(w_{\nu}^{(i)})\in\overline{E}^{(i)}\ (1\leqslant\nu\leqslant f_i)$ 是 $\overline{F}$-线性无关的, 而 $\varphi^{(i)}(\pi_{\mu}^{(i)})\ (0\leqslant\mu\leqslant e_i-1)$ 为 $\varphi^{(i)}(E^{\times})$ 对于 $\varphi(F^{\times})$ 的陪集完全代表系. 我们只需证明下列 $\sum_{i=1}^{g}e_if_i$ 个元素
\[
  \{w_{\nu}^{(i)}\pi_{\mu}^{(i)}\mid 1\leqslant i\leqslant g,\ 1\leqslant\nu\leqslant f_i,\ 0\leqslant\mu\leqslant e_i-1\}
\]
是 $F$-线性无关的. 若不然, 则有不全为 $0$ 的一组元素 $a_{\nu\mu}^{(i)}\in F$, 使得
\[
  \sum_{i,\nu,\mu}a_{\nu\mu}^{(i)}w_{\nu}^{(i)}\pi_{\mu}^{(i)}=0.
\]
不妨设
\[
  \varphi(a_{10}^{(1)})=\max_{i,\nu,\mu}\{\varphi(a_{\nu\mu}^{(i)})\},
\]
则 $\varphi(a_{10}^{(1)})\neq 0$. 由引理 \ref{lem:ch7-10} 的证明可知
\[
  M=\varphi^{(1)}\left(\sum_{\nu,\mu}a_{\nu\mu}^{(1)}w_{\nu}^{(1)}\pi_{\mu}^{(1)}\right)
  =\max_{\nu,\mu}\{\varphi^{(1)}(a_{\nu\mu}^{(1)}\pi_{\mu}^{(1)})\}
  \geqslant\varphi^{(1)}(a_{10}^{(1)}\pi_0^{(1)})>0.
\]
我们要适当选取 $w_{\nu}^{(i)}\ (2\leqslant i\leqslant g,\ 1\leqslant\nu\leqslant f_i)$, 使得对任意 $\nu$ 和 $\mu$ 均有
\begin{equation}\tag{*}
  \varphi^{(1)}\left(a_{\nu\mu}^{(i)}w_{\nu}^{(i)}\pi_{\mu}^{(i)}\right)<M\qquad(2\leqslant i\leqslant g).
\end{equation}
这时,
\[
  0=\varphi^{(1)}(0)=\varphi^{(1)}\left(\sum_{i,\nu,\mu}a_{\nu\mu}^{(i)}w_{\nu}^{(i)}\pi_{\mu}^{(i)}\right)
  =\varphi^{(1)}\left(\sum_{\nu,\mu}a_{\nu\mu}^{(1)}w_{\nu}^{(1)}\pi_{\mu}^{(1)}\right)=M>0,
\]
这个矛盾即证明了定理. 为了选取满足 (*) 式的 $w_{\nu}^{(i)}$, 要用逼近定理 \ref{thm:7.4}. 令
\[
  N=\min_{\substack{i,\mu\\ i\neq 1}}\varphi^{(1)}(\pi_0^{(1)})/\varphi^{(1)}(\pi_{\mu}^{(i)}),
\]
我们只要使 $\{w_{\nu}^{(i)}\mid 2\leqslant i\leqslant g,\ 1\leqslant\nu\leqslant f_i\}$ 满足
\[
  \varphi^{(1)}(w_{\nu}^{(i)})<N\qquad(2\leqslant i\leqslant g,\ 1\leqslant\nu\leqslant f_i)
\]
即可. 因为这时对 $i\geqslant 2$,
\begin{align*}
  \varphi^{(1)}\left(a_{\nu\mu}^{(i)}w_{\nu}^{(i)}\pi_{\mu}^{(i)}\right)
  &<N\cdot\varphi^{(1)}(\pi_{\mu}^{(i)}/\pi_0^{(1)})\cdot\varphi^{(1)}(a_{10}^{(1)}\pi_0^{(1)})\\
  &\leqslant N\cdot N^{-1}\cdot M=M.
\end{align*}
但是, 若原先选取的为 $\{\widetilde{w}_{\nu}^{(i)}\}$, 则对每个固定的 $i\geqslant 2$ 和 $\nu$, 由逼近定理可知存在 $w_{\nu}^{(i)}\in E$, 使得
\[
  \varphi^{(i)}\left(\widetilde{w}_{\nu}^{(i)}-w_{\nu}^{(i)}\right)<1,\qquad
  \varphi^{(1)}\left(w_{\nu}^{(i)}\right)<N\qquad(2\leqslant i\leqslant g).
\]
由第 $1$ 式和 $\psi^{(i)}(w_{\nu}^{(i)})=\psi^{(i)}(\widetilde{w}_{\nu}^{(i)})$. 从而新的 $\{w_{\nu}^{(i)}\}$ 即为所求.
\end{proof}
